\documentclass[10pt,a4paper]{amsart}
\usepackage[margin=19mm]{geometry}
\usepackage{amsmath,amssymb,mathtools}
\usepackage[scr=rsfs,cal=euler]{mathalfa}
\usepackage{xfrac}
\usepackage[unicode,hidelinks,bookmarks=false]{hyperref}
\newcommand{\ie}{i.e.\ }
\newcommand{\cf}{cf.\ }
\newcommand{\resp}{respectively\ }
\newcommand{\Subsection}{\subsection}
\providecommand{\nobreakdash}{\nobreak\hbox{-}}
\setlength{\emergencystretch}{2em}
\multlinegap0pt
\usepackage{amssymb}
\usepackage{xcolor}
\usepackage{tikz}
\usepackage{enumerate}
\newtheorem{lemma}{Lemma}
\title{A note on the Brill-Noether loci of small codimension in moduli space of stable bundles}
\author{Pritthijit Biswas, Jaya NN Iyer}
\date{Source: arXiv:2505.15749v2 (CC BY 4.0)}
\begin{document}
\maketitle

\noindent\emph{Source and licence.} A note on the Brill-Noether loci of small codimension in moduli space of stable bundles. Version arXiv:2505.15749v2 (CC BY 4.0), distributed by arXiv under the Creative Commons Attribution 4.0 International licence, http://creativecommons.org/licenses/by/4.0/, as recorded in the arXiv licence metadata for this exact version and shown on the abstract page https://arxiv.org/abs/2505.15749v2. Full licence notice: \texttt{licenses/arxiv-2505.15749-CC-BY-4.0.txt}.\par\smallskip

\noindent\textit{Editorial scope.} Counted unit: the article's lemma 4.33 (source0.tex lines 396--399). The article's complete original proof of it is delivered with it (source0.tex lines 400--416, one \texttt{proof} environment of 17 lines). No mathematics is written, repaired or re-derived: the statement and the proof are the article's own lines, in the article's own notation, and no retained line is substituted or re-flowed.  No further source-local context is needed: every cross-reference of every retained line resolves inside the two delivered spans.  Nothing was omitted from any retained span, so the package carries no omission record.  No retained line contains a citation, so the wrapper carries no bibliography and no bibliography entry.  No retained line contains a figure environment, an image-inclusion command or drawing code, so no image is delivered and the package declares no asset dependency.  Editorial formatting only, in addition: an AMS \texttt{amsart} preamble that restates the article's own declarations and macros used by the retained lines, namely the environments \texttt{lemma} on separate counters, together with the standard packages those lines need.  The retained environments print in this document as Lemma 1 in order of appearance, whereas the article prints the same items with the numbering of its own full document.  The complete native source of the article as distributed in the arXiv e-print for this exact version is bundled unchanged as \texttt{sources/178535/source0.tex}.
\begin{lemma}\label{4.33}
Let $X\subset \mathbb{P}^n$ be any hypersurface of degree $d$.
The family of lines $\mathcal{L}(S)$ surjects onto $X$ if $\mbox{dim}~S\geq d$.
\end{lemma}

\begin{proof} Let $p=(p_i)\in X$ be a general point.
Consider the polar hypersurfaces of $X$ with respect to $p$:
\begin{eqnarray*}
P_1(x,p) & : &\sum_{i=1}^n\frac{\partial{F}}{dx_i}(x).p_i=0 \\
P_2(x,p) & : & \sum\frac{\partial{F}}{\partial{x_i}\partial{x_j}}(x).p_ip_j=0 \\
&& \vdots \\
P_{d-1}(x,p)& : &\sum \frac{\partial{F}}{\partial{x_{i_1}}\partial{x_{i_2}}...\partial{x_{i_{(d-1)}}}}(x).p_{i_1}p_{i_2}...p_{i_{(d-1)}}=0.
\end{eqnarray*}

Then clearly the variety $\mathbb{P}(p)$ defined by the equations 
$P_1=P_2=...=P_{d-1}=0$ 
has a non-empty intersection $Z$ with $S$. This means that for $z\in Z$ the 
line 
$l_{zp}$ joining $p$ and $z$ has intersection $l_{zp}.X=d-1$ if $l_{zp}\not\subset X$. 
 In particular since dimension of any irreducible component of $Z$ is at least $1$ there is a $z\in Z$ and $z\neq p$
such that $p\in \mathcal{V}(z)$, i.e., $l_{zp}\subset X$. 
\end{proof}

\end{document}
